% ------------------------------------------------------------------
% Baseline scenario without multivariate selection (referee request).
% Suggested placement: in Section "Results", after the discussion of
% Table~\ref{tab:results_xsec}; or as a short subsection at the end of
% the statistical-interpretation section.
%
% NUMBERS: obtained with the identical event selections, samples and
% fit machinery as the nominal results (SDCC reproduction pipeline,
% whose nominal precisions agree with Table~\ref{tab:results_xsec}
% within 10--20\%). Swap in official-pipeline numbers if regenerated.
% ------------------------------------------------------------------

To assess the dependence of the results on the multivariate selection, a
baseline scenario is evaluated in which no BDT information is used at any
stage. The event selection is restricted to the cut-based requirements of
Section~\ref{sec:EventSelection}, and the cross section is extracted for every
final state from a one-dimensional fit to the recoil-mass distribution in
the full $100 < m_{\text{recoil}} < 150\,\GeV$ range, in a single region. In this
configuration the background normalizations are free parameters of the
fit, determined in situ from the recoil-mass sidebands, in close analogy
to the background-estimation strategies employed in the LEP
Higgs-boson searches; only the background \emph{shapes} are taken from
simulation. The resulting uncertainties are summarized in
Table~\ref{tab:results_baseline}.

The comparison with Table~\ref{tab:results_xsec} quantifies the impact of
the multivariate selection. At \sqrtsZH the leptonic precisions degrade by
factors of 1.3--1.5, reflecting the fact that the recoil-mass peak itself
carries most of the discriminating power at this energy, while at
\sqrtsTop the degradation reaches a factor of about 2.4 per leptonic
channel, as the recoil distribution is broadened by initial-state
radiation and the larger beam-energy spread. In the hadronic final state
the baseline precision degrades by a factor of about three at \sqrtsZH:
beyond the loss in purity, the removal of the low-BDT control region
leaves the individual background normalizations nearly degenerate when
constrained by the recoil shape alone, and the associated normalization
uncertainties dominate the baseline result. The combined precision
degrades from $\pm 0.31\%$ ($\pm 0.52\%$) to $\pm 0.64\%$ ($\pm 1.31\%$)
at \sqrtsZH (\sqrtsTop). The baseline scenario thus validates the
in-situ background determination underlying the nominal analysis, while
demonstrating that the multivariate selection improves the final
precision by a factor of about two at both center-of-mass energies
without introducing a dependence on the Higgs-boson decay composition
beyond that quantified in Section~\ref{sec:Model-independent}.

\begin{table}[ht!]
\renewcommand{\arraystretch}{1.1}
\centering
\caption{Relative uncertainties at 68\% confidence level on the \ZH cross
section (in percent) for the baseline scenario without multivariate
selection: cut-based selection and one-dimensional recoil-mass fit with
background normalizations determined from the sidebands, shown for the
individual final states and their combination.}
\vspace{0.5em}
\label{tab:results_baseline}
{\small
\begin{tabular}{lcc}
\toprule
\textbf{Final state} & \textbf{\sqrtsZH} & \textbf{\sqrtsTop} \\
\midrule
\multicolumn{3}{l}{\textbf{Leptonic channels}} \\
\ZeeH   & $\pm 1.15$ & $\pm 5.20$ \\
\ZmumuH & $\pm 0.88$ & $\pm 4.59$ \\
\midrule
\ZllH   & $\pm 0.71$ & $\pm 3.53$ \\
\midrule
\multicolumn{3}{l}{\textbf{Hadronic channel}} \\
\ZqqH   & $\pm 1.45$ & $\pm 1.40$ \\
\midrule
\multicolumn{3}{l}{\textbf{Combined result}} \\
\ZH     & $\pm 0.64$ & $\pm 1.31$ \\
\bottomrule
\end{tabular}
}
\end{table}
